\section{Finite completion and explicit competitors}
\label{sec:finite}

The preceding analysis reduces each remaining statement to an explicitly
bounded set of rational positive-definiteness tests. We give the complete
ranges and the separating inequalities, including the cases where a
stronger proposed cap is false.

\subsection{The one-cell family and its uniform ceiling}

\begin{proof}[Proof of Theorem~\ref{thm:one-cell}]
For $k\ge25$, the single legal cell has length $8k+2\ge202$, so
Theorem~\ref{thm:cells} at $r=1$, $\alpha=+1$, and $c=\cone$ applies.
The remaining $k=1,\ldots,24$ are exactly the orders
\[
 10,18,26,\ldots,194.
\]
For each full signed graph the rational verifier proves
$\cone I-A_{8k+2}(w_k,+1)^2\succ0$ by positive scalar Schur pivots.
The $10\times10$ matrix is constructed and checked directly, without
assuming disjoint ordinary and wrap block locations. The first analytic
order is $202$, so no admissible size is missing.

At $c_-=7905369/1000000$ and order $202$, independent exact elimination
of $c_-I-A_{202}(w_{25},+1)^2$ has $196$ positive scalar interior pivots.
The retained six-coordinate core then has a strictly negative pivot after
a positive scalar-pivot prefix. The recorded rational pivot is strictly
negative, not merely nonpositive or a failed numerical test.
Schur congruence gives a vector with negative quadratic form, and hence
\[
 \rho(A_{202}(w_{25},+1))^2>c_-.
\]
This member supplies the lower endpoint in \eqref{eq:supremum}; the
uniform strict upper bounds supply the non-strict upper endpoint.
\end{proof}

\subsection{Two and three nearly equal legal cells}

For $r=2$ and $k\ge6$, the parameters
$\lfloor k/2\rfloor,\lceil k/2\rceil$ sum to $k$.
When $k\ge26$, each is at least $13$, so the cell bound $\czero$ applies.
Exactly twenty finite cases remain:
\begin{equation}
 k=6,\ldots,25,\qquad N=52,60,\ldots,204.
 \label{eq:R4finite}
\end{equation}
All use negative holonomy. For even $k$ their cell lengths are equal.
For odd $k$ the two lengths are $N/2-4$ and $N/2+4$, which changes the
placement from the older equally spaced gap construction. Each matrix in
\eqref{eq:R4finite} is rebuilt from its own word and passes exact rational
LDL positivity at cap $\czero$. The first analytic order is $212$.

For $r=3$, write $k=3q+s$, $0\le s<3$. The three parameters
\[
 \left\lfloor\frac{k}{3}\right\rfloor,
 \left\lfloor\frac{k+1}{3}\right\rfloor,
 \left\lfloor\frac{k+2}{3}\right\rfloor
\]
are $3-s$ copies of $q$ followed by $s$ copies of $q+1$.
They sum to $k$, are at least two for $k\ge6$, and differ by at most one.
For $k\ge39$ they are all at least thirteen, so Theorem~\ref{thm:cells}
applies with positive holonomy. The remaining thirty-three cases are
\begin{equation}
 k=6,\ldots,38,\qquad N=54,62,\ldots,310.
 \label{eq:R6finite}
\end{equation}
Each full signed graph passes exact rational LDL at cap $\czero$.
The first analytic order is $318$, immediately after the final admissible
finite order. When $3\nmid k$, the cells are nearly equal rather than
identical, and no certificate for another gap placement is substituted.

For both lists, the verification forms $(198/25)I-A^2$ or its
denominator-cleared equivalent $198I-25A^2$ directly from integer
two-edge walks. A separate implementation uses the normalized matrix and
a different scalar elimination order. For the three-cell list,
it also reconstructs the triangle signs independently from the
quadrilateral-gap word. Every required pivot in each list is strictly
positive. These are the only finite ranges left by the minimum-cell-length
reduction for the stated residue-four and residue-six families.

\subsection{The twisted comparison and the all-even witness range}

\begin{proof}[Proof of Theorem~\ref{thm:witnesses}]
For residue two, Theorem~\ref{thm:one-cell} gives the stronger cap
$\cone<\czero$ at every $k\ge6$. The two- and three-cell arguments above
give the cap $\czero$ for the other two residues.
For $x>0$, $\cos x>1-x^2/2$, and $\pi^2<10$. Therefore
\begin{equation}
 \tw(N)^2>8-\frac{20\pi^2}{N^2}>8-\frac{200}{N^2}.
 \label{eq:benchmark-lower}
\end{equation}
At $N\ge50$, the last expression is at least $198/25$.
Thus every nonzero residue witness satisfies
\eqref{eq:nonzero-witnesses}. In particular, the exact endpoint gaps for
the other two residues are
\[
 \frac{2679}{338}-\frac{198}{25}=\frac{51}{8450}>0,
 \qquad
 \frac{5782}{729}-\frac{198}{25}=\frac{208}{18225}>0,
\]
coming from $N=52$ and $N=54$ respectively.

For $N=8m\ge32$, Theorem~\ref{thm:period8} gives a witness with squared
radius at most $\etaedge$. The rational comparison
\begin{equation}
 \etaedge<\frac{999}{128}=8-\frac{200}{32^2}
 \label{eq:period8-separator}
\end{equation}
follows by squaring positive quantities: the intermediate bound is
\[
 \frac{(999/128-4)^2-10}{2}=\frac{73329}{32768}>\sqrt5,
 \quad 73329^2-5\cdot32768^2=8433121>0.
\]
Equations~\eqref{eq:benchmark-lower} and
\eqref{eq:period8-separator} prove the strict comparison at every such
order. The four residue classes now cover $32,40$ and every even $N\ge48$.
\end{proof}

\subsection{A stronger balanced subsequence and exact obstructions}

Equal cells admit one additional small-order completion.
\begin{proposition}\label{prop:balanced-sharp}
For every $j\ge1$,
\[
 \rho(A_{16j+4}(w_j\|w_j,-1))^2<\cone.
\]
For $j\ge3$ this value is strictly below the twisted benchmark.
\end{proposition}
\begin{proof}
Corollary~\ref{cor:phase} applies for $j\ge25$, with the two finite
phases $z=\pm\mathrm i$. The remaining $j=1,\ldots,24$, at orders
$20,36,\ldots,388$, have exact full-graph positive-definiteness certificates
for $790537I-100000A^2$. Both primary and independent rational replays
cover the complete list. Their smallest case uses a cell of length ten
and is checked directly, including the coincident couplings.
The first analytic order is $404$. For $j\ge3$, $N\ge52$, so
\eqref{eq:benchmark-lower} and $\cone<\czero$ give the comparison.
\end{proof}

The sharper cap cannot replace $\czero$ in the full residue-four theorem.
For the negative-holonomy nearly balanced words, exact rational elimination
gives
\begin{align*}
 \rho(A_{60}(w_3\|w_4,-1))^2&>\cone,\\
 \rho(A_{76}(w_4\|w_5,-1))^2&>\cone.
\end{align*}
In each case the certificate records a strictly negative pivot with its
complete positive predecessor prefix, proving a negative direction for
$\cone I-A^2$. Thus the obstruction is exact and applies to the actual
changed construction used here.
There is also a distinct obstruction for the older balanced $N=60$ gap
word $[6,4^6,6,4^6]$ with negative holonomy. Its half-length is thirty;
its triangle signs change sign under that half-shift rather than repeating
a legal $w_j$. It therefore lies outside
Proposition~\ref{prop:balanced-sharp}, and its own exact negative pivot
confirms that dropping the even-$k$ condition would be false.
None of these obstructions contradicts the $\czero$ bound or gives an
unrestricted lower bound over all signings.
